% Part: incompleteness
% Chapter: incompleteness-provability
% Section: rosser-thm

\documentclass[../../../include/open-logic-section]{subfiles}

\begin{document}

\olfileid{inc}{inp}{ros}

\olsection{د روسر قضيه}

ايا د ګوډل ثبوت داسې بدلولے شو چې د «$\omega$-سازګارې» پر ځاے
يوازې «سازګاره» تيوري فرض کړو او پياوړې پايله واخلو؟ ځواب «هو» دے؛
د روسر موندلې چاره دا کار کوي۔ د روسر چاره دا ده چې د
$\OProv[\Th{T}](y)$ پر ځاے د !!{derivability} يو «بدل شوے»
پريديکات~$\ORProv_T(y)$ وکاروو۔

\begin{thm}
\ollabel{thm:rosser}
فرض کړئ $\Th{T}$ هره سازګاره او !!{axiomatizable} تيوري ده چې
$\Th{Q}$ غځوي۔ نو~$\Th{T}$ بشپړه نۀ ده۔
\end{thm}

\begin{proof}
ياد ولرئ چې $\OProv[\Th{T}](y)$ د
$\lexists[x][\OPrf[\Th{T}](x,
  y)]$ په توګه تعريف شوے دے۔ دلته~$\OPrf[\Th{T}](x, y)$ هغه
د پرېکړې وړ اړيکه تمثيلوي چې هله او يوازې هله برقرارېږي چې~$x$
د ګوډل شمېرې~$y$ لرونکې !!{sentence} د !!a{derivation} ګوډل
شمېره وي۔ هغه اړيکه هم د پرېکړې وړ ده چې د~$x$ او~$y$ تر منځ
هله برقرارېږي چې~$x$ د ګوډل شمېرې~$y$ لرونکې جملې د
\emph{ردونې} ګوډل شمېره وي۔ $\fn{not}(x)$ دې هغه ابتدايي
بازګشتي تابعه وي چې داسې کار کوي: که~$x$ د يوه فورمول~$!A$ کوډ
وي، نو~$\fn{not}(x)$ د~$\lnot
!A$ کوډ وي۔ په دې توګه $\Refut[\Th{T}](x, y)$ هله او يوازې هله
برقرارېږي چې $\Prf[\Th{T}](x, \fn{not}(y))$ برقرار وي۔
$\ORefut[\Th{T}](x, y)$ دې دا اړيکه تمثيل کړي۔ نو که
$\Th{T} \Proves \lnot !A$ وي او~$\delta$ ئې اړوند !!{derivation}
وي، $\Th{Q} \Proves
\ORefut[\Th{T}](\gn{\delta}, \gn{!A})$۔ موږ
$\ORProv[\Th{T}](y)$ داسې تعريفوو:
\[
\lexists[x][(\OPrf[\Th{T}](x,y) \land \lforall[z][(z < x \lif \lnot
  \ORefut[\Th{T}](z,y))])].
\]
په تقريبي ډول~$\ORProv[\Th{T}](y)$ وايي: «په~$\Th{T}$ کښې د~$y$
ثبوت شته، او د~$y$ تر دې لنډه ردونه نشته»۔ که~$\Th{T}$ سازګاره
فرض کړو، نو~$\ORProv[\Th{T}](y)$ د هماغو عددونو دپاره رښتيا دے
چې~$\OProv[\Th{T}](y)$ ئې دپاره رښتيا دے؛ خو په~$\Th{T}$ کښې د
\emph{اثبات‌وړتيا} له پلوه (او اوس پوهېږو چې د رښتياوالي او
اثبات‌وړتيا تر منځ توپير شته!) دواړه بېل خاصيتونه لري۔ که
$\Th{T}$ \emph{نا}سازګاره وي، نو دواړه د هماغو عددونو دپاره
\emph{نۀ} برقرارېږي!{} ($\ORProv[\Th{T}](y)$ ته ډېر ځله داسې
لوستل کېږي: «$y$ د روسر له مخې اثبات‌وړ دے»۔ خو لکه چې مو همدا
اوس وليدل، د روسر اثبات‌وړتيا د اثبات‌وړتيا کوم ځانګړے ډول نۀ
دے: په ناسازګارو تيوريو کښې داسې !!{sentence}s شته چې اثبات‌وړ
دي، خو د روسر له مخې اثبات‌وړ نۀ دي۔ دا نوم لوستونکے ګډوډولے شي؛
د ګډوډۍ د مخنيوي دپاره ورته يو بېل مصنوعي نوم، لکه «$y$ شمووړ دے»، هم
ويلے شئ۔)

د ثابت ټکي د لمې له مخې داسې فورمول~$!R_\Th{T}$ شته چې
\begin{equation}
  \Th{Q} \Proves !R_\Th{T} \liff \lnot \ORProv[\Th{T}](\gn{!R_\Th{T}}).
  \ollabel{RT}
\end{equation}
د~\olref[1in]{thm:first-incompleteness} د ثبوت برعکس، دلته دعوه
دا ده چې که~$\Th{T}$ سازګاره وي، نو~$\Th{T}$، $!R_\Th{T}$
نۀ !!{derive} کوي، او~$\Th{T}$، $\lnot !R_\Th{T}$ هم نۀ !!{derive} کوي۔
(په نورو ټکو، د $\omega$-سازګارۍ فرض ته اړتيا نشته۔)

لومړے ښيو چې $\Th{T} \Proves/ !R_{\Th{T}}$۔ که ئې ثابت کړے
واے، نو له~$T$ څخه به د~$!R_{\Th{T}}$ يو !!a{derivation} و؛ د
هغې ګوډل شمېره دې~$n$ وي۔ بيا
$\Th{Q} \Proves \OPrf[\Th{T}](\num{n}, \gn{!R_{\Th{T}}})$، ځکه چې
$\OPrf[\Th{T}]$ په~$\Th{Q}$ کښې $\Prf[\Th{T}]$ تمثيلوي۔ همدا
راز، د هر~$k < n$ دپاره~$k$ د~$\lnot !R_{\Th{T}}$ د يوه
!!a{derivation} ګوډل شمېره نۀ ده، ځکه~$\Th{T}$ سازګاره ده۔ نو د
هر~$k < n$ دپاره
$\Th{Q} \Proves \lnot
\ORefut[\Th{T}](\num{k}, \gn{!R_{\Th{T}}})$۔ د
\olref[req][min]{lem:less-nsucc} له مخې
$\Th{Q} \Proves \lforall[z][(z < \num{n} \lif \lnot \ORefut[\Th{T}](z,
  \gn{!R_{\Th{T}}}))]$۔ نو
\[
\Th{Q} \Proves \lexists[x][(\OPrf[\Th{T}](x,\gn{!R_{\Th{T}}}) \land \lforall[z][(z
    < x \lif \lnot \ORefut[\Th{T}](z,\gn{!R_{\Th{T}}}))])],
\]
خو دا هماغه~$\ORProv[\Th{T}](\gn{!R_{\Th{T}}})$ دے۔ د~\olref{RT}
له مخې $\Th{Q}
\Proves \lnot !R_{\Th{T}}$۔ ځکه چې~$\Th{T}$، $\Th{Q}$ غځوي،
نو $\Th{T}
\Proves \lnot !R_{\Th{T}}$ هم دے۔ موږ فرض کړې وه چې
$\Th{T} \Proves !R_{\Th{T}}$؛ نو~$\Th{T}$ به ناسازګاره شي، چې د
قضيې له فرض سره ټکر لري۔

اوس ښيو چې $\Th{T} \Proves/ \lnot !R_{\Th{T}}$۔ بيا فرض کړئ چې
دا نفي ثابتوي او~$n$ د~$\lnot !R_{\Th{T}}$ د يوه
!!a{derivation} ګوډل شمېره ده۔ نو
$\Refut[\Th{T}](n, \Gn{!R_{\Th{T}}})$ برقرارېږي۔ ځکه چې
$\ORefut[\Th{T}]$ په~$\Th{Q}$ کښې $\Refut[\Th{T}]$ تمثيلوي،
نو $\Th{Q} \Proves
\ORefut[\Th{T}](\num{n}, \gn{!R_{\Th{T}}})$۔ بيا به وښيو چې
$\Th{T}$ ناسازګاره کېږي، ځکه~$!R_{\Th{T}}$ به هم !!{derive} کوي۔
څرنګه چې
\begin{align*}
\Th{Q} & \Proves !R_{\Th{T}} \liff \lnot \ORProv[\Th{T}](\gn{!R_{\Th{T}}}),
\intertext{او ځکه چې $\Th{T}$، $\Th{Q}$ غځوي، بس دا وښيو چې}
\Th{Q} & \Proves \lnot \ORProv[\Th{T}](\gn{!R_{\Th{T}}}).
\end{align*}
!!{sentence}~$\lnot
\ORProv[\Th{T}](\gn{!R_{\Th{T}}})$، يعنې
\begin{align*}
  \lnot & \lexists[x][(\OPrf[\Th{T}](x,\gn{!R_{\Th{T}}}) \land
    \lforall[z][(z < x \lif \lnot \ORefut[\Th{T}](z,\gn{!R_{\Th{T}}}))])],
  \intertext{په منطقي ډول له دې سره معادله ده:}
  & \lforall[x][(\OPrf[\Th{T}](x,\gn{!R_{\Th{T}}}) \lif
    \lexists[z][(z < x \land \ORefut[\Th{T}](z,\gn{!R_{\Th{T}}}))])].
\end{align*}
موږ د منطق او په~$\Th{Q}$ کښې د !!{derive}s کېدونکو حقايقو په
کارولو غير صوري استدلال کوو۔ فرض کړئ~$x$ خپل‌سرے دے او
$\OPrf[\Th{T}](x,
\gn{!R_{\Th{T}}})$ برقرارېږي۔ موږ لا پوهېږو چې
$\Th{T} \Proves/ !R_{\Th{T}}$، نو د هر~$k$ دپاره
$\Th{Q} \Proves \lnot \OPrf[\Th{T}](\num{k}, \gn{!R_{\Th{T}}})$۔
له دې د هر~$k$ دپاره $\eq/[x][\num{k}]$ راځي۔ په ځانګړي ډول (الف)
$\eq/[x][\num{n}]$ لرو۔ همدا راز
$\lnot(\eq[x][\num{0}]
\lor \eq[x][\num{1}] \lor \dots \lor \eq[x][\num{n-1}])$ لرو،
نو د~\olref[req][min]{lem:less-nsucc} له مخې (ب)
$\lnot(x < \num{n})$۔ د~\olref[req][min]{lem:trichotomy} له مخې
$\num{n} < x$۔ ځکه چې
$\Th{Q} \Proves
\ORefut[\Th{T}](\num{n}, \gn{!R_{\Th{T}}})$، لرو چې
$\num{n} < x \land
\ORefut[\Th{T}](\num{n}, \gn{!R_{\Th{T}}})$؛ له دې څخه
$\lexists[z][(z < x
\land \ORefut[\Th{T}](z,\gn{!R_{\Th{T}}}))]$ راځي۔ ځکه چې~$x$
خپل‌سرے و، اړينه پايله دا ده:
\[
\lforall[x][(\OPrf[\Th{T}](x,\gn{!R_{\Th{T}}}) \lif
  \lexists[z][(z < x \land \ORefut[\Th{T}](z,\gn{!R_{\Th{T}}}))])].
\]
\end{proof}

\begin{prob}
د طبيعي عددونو دوه سټونه~$A$ او~$B$ هغه وخت
\emph{په محاسبوي ډول نه بېلېدونکي} بلل کېږي چې هېڅ !!{decidable}
سټ~$X$ داسې نۀ وي چې $A \subseteq X$ او
$B \subseteq \Complement{X}$ وي ($\Complement{X}$ د~$X$
متمم، يعنې~$\Nat \setminus X$، دے)۔ فرض کړئ~$\Th{T}$ د~$\Th{Q}$
سازګاره او !!{axiomatizable} غځونه ده۔ $A$ دې په~$\Th{T}$ کښې د
اثبات‌وړو !!{sentence}s د ګوډل شمېرو سټ وي، او~$B$ دې په~$\Th{T}$
کښې د ردېدونکو جملو د ګوډل شمېرو سټ وي۔ ثابت کړئ چې~$A$ او~$B$
په محاسبوي ډول نه بېلېدونکي دي۔
\end{prob}

\end{document}
